\input{preamble/preamble.tex}

\geometry{margin=0.85in}
\AtBeginDocument{\small}

\newcommand{\ExamNameField}{\noindent\textbf{Name:}\ \rule{0.7\linewidth}{0.4pt}\par\medskip}

\begin{document}
	\ExamTitleBlock{11th grade}{Term 2 Learning Evidence 2: C6 Sample Size Activity}{}
	
	\section*{1. Renewable Energy Subsidy Approval}
		

	A policy team surveys \textbf{200 households} across a region to estimate the share that approve a proposed renewable energy subsidy. In the sample, \textbf{128 households} say they approve.
	
	Construct and interpret a \textbf{90\% confidence interval} for the population proportion of households that approve the subsidy.
	
	\subsubsection*{(C7) Constructs the confidence interval for a proportion}
	To construct the confidence interval for a proportion, we identify \(x\) and \(n\) and compute \(\hat{p}\), check normal approximation conditions \(n\hat{p} \ge 10\) and \(n(1-\hat{p}) \ge 10\), choose the correct \(z^*\) for the confidence level, compute \(SE\) and the margin of error, construct and report the confidence interval, and interpret the interval in context.
	
	\textbf{Step 1: Sample proportion.}
	\[
	\hat{p} = \frac{x}{n} = \frac{128}{200} = 0.64
	\]
	
	\textbf{Step 2: Large-sample conditions.}
	\[
	n\hat{p} = 200(0.64) = 128 \ge 10, \qquad n(1-\hat{p}) = 200(0.36) = 72 \ge 10
	\]
	
	\textbf{Step 3: Standard error.}
	\[
	SE = \sqrt{\frac{\hat{p}(1-\hat{p})}{n}} = \sqrt{\frac{0.64(0.36)}{200}} \approx 0.0339
	\]
	
	\textbf{Step 4: Margin of error (90\% confidence, \(z^* = 1.645\)).}
	\[
	E = z^* \cdot SE = 1.645(0.0339) \approx 0.0558
	\]
	
	\textbf{Step 5: Confidence interval.}
	\[
	\hat{p} \pm E = 0.64 \pm 0.0558 \Rightarrow [0.584,\ 0.696]
	\]
	
	\textbf{Interpretation.}
	We are 90\% confident that between about 58.4\% and 69.6\% of all households in the region approve the renewable energy subsidy.
	
	\section*{2. Export Shipment Delays}
	
	An export council reviews \textbf{150 shipments} and finds that \textbf{39} were delayed at least one day. The council wants a proportion estimate for the delay rate in the current trade season.
	
	Construct and interpret a \textbf{95\% confidence interval} for the population proportion of delayed shipments.
	
	\subsubsection*{(C7) Constructs the confidence interval for a proportion}
	To construct the confidence interval for a proportion, we identify \(x\) and \(n\) and compute \(\hat{p}\), check normal approximation conditions \(n\hat{p} \ge 10\) and \(n(1-\hat{p}) \ge 10\), choose the correct \(z^*\) for the confidence level, compute \(SE\) and the margin of error, construct and report the confidence interval, and interpret the interval in context.
	
	\textbf{Step 1: Sample proportion.}
	\[
	\hat{p} = \frac{x}{n} = \frac{39}{150} = 0.26
	\]
	
	\textbf{Step 2: Large-sample conditions.}
	\[
	n\hat{p} = 150(0.26) = 39 \ge 10, \qquad n(1-\hat{p}) = 150(0.74) = 111 \ge 10
	\]
	
	\textbf{Step 3: Standard error.}
	\[
	SE = \sqrt{\frac{\hat{p}(1-\hat{p})}{n}} = \sqrt{\frac{0.26(0.74)}{150}} \approx 0.0358
	\]
	
	\textbf{Step 4: Margin of error (95\% confidence, \(z^* = 1.96\)).}
	\[
	E = z^* \cdot SE = 1.96(0.0358) \approx 0.0702
	\]
	
	\textbf{Step 5: Confidence interval.}
	\[
	\hat{p} \pm E = 0.26 \pm 0.0702 \Rightarrow [0.190,\ 0.330]
	\]
	
	\textbf{Interpretation.}
	We are 95\% confident that between about 19.0\% and 33.0\% of all export shipments are delayed at least one day this season.
	
	\section*{3. Microloan Repayment Defaults}
	
	A microfinance program tracks \textbf{400 loans} issued to small businesses. Within the first year, \textbf{50 loans} defaulted.
	
	Construct and interpret a \textbf{98\% confidence interval} for the population proportion of loans that default within the first year.
	
	\subsubsection*{(C7) Constructs the confidence interval for a proportion}
	To construct the confidence interval for a proportion, we identify \(x\) and \(n\) and compute \(\hat{p}\), check normal approximation conditions \(n\hat{p} \ge 10\) and \(n(1-\hat{p}) \ge 10\), choose the correct \(z^*\) for the confidence level, compute \(SE\) and the margin of error, construct and report the confidence interval, and interpret the interval in context.
	
	\textbf{Step 1: Sample proportion.}
	\[
	\hat{p} = \frac{x}{n} = \frac{50}{400} = 0.125
	\]
	
	\textbf{Step 2: Large-sample conditions.}
	\[
	n\hat{p} = 400(0.125) = 50 \ge 10, \qquad n(1-\hat{p}) = 400(0.875) = 350 \ge 10
	\]
	
	\textbf{Step 3: Standard error.}
	\[
	SE = \sqrt{\frac{\hat{p}(1-\hat{p})}{n}} = \sqrt{\frac{0.125(0.875)}{400}} \approx 0.0165
	\]
	
	\textbf{Step 4: Margin of error (98\% confidence, \(z^* = 2.326\)).}
	\[
	E = z^* \cdot SE = 2.326(0.0165) \approx 0.0385
	\]
	
	\textbf{Step 5: Confidence interval.}
	\[
	\hat{p} \pm E = 0.125 \pm 0.0385 \Rightarrow [0.087,\ 0.164]
	\]
	
	\textbf{Interpretation.}
	We are 98\% confident that between about 8.7\% and 16.4\% of loans in the full portfolio default within the first year.
	
	\section*{4. Inflation Expectations Survey}
	
	A national statistics office surveys \textbf{120 households} about whether they expect inflation to rise over the next year. \textbf{78 households} report expecting higher inflation.
	
	Construct and interpret a \textbf{99\% confidence interval} for the population proportion of households expecting higher inflation.
	
	\subsubsection*{(C7) Constructs the confidence interval for a proportion}
	To construct the confidence interval for a proportion, we identify \(x\) and \(n\) and compute \(\hat{p}\), check normal approximation conditions \(n\hat{p} \ge 10\) and \(n(1-\hat{p}) \ge 10\), choose the correct \(z^*\) for the confidence level, compute \(SE\) and the margin of error, construct and report the confidence interval, and interpret the interval in context.
	
	\textbf{Step 1: Sample proportion.}
	\[
	\hat{p} = \frac{x}{n} = \frac{78}{120} = 0.65
	\]
	
	\textbf{Step 2: Large-sample conditions.}
	\[
	n\hat{p} = 120(0.65) = 78 \ge 10, \qquad n(1-\hat{p}) = 120(0.35) = 42 \ge 10
	\]
	
	\textbf{Step 3: Standard error.}
	\[
	SE = \sqrt{\frac{\hat{p}(1-\hat{p})}{n}} = \sqrt{\frac{0.65(0.35)}{120}} \approx 0.0435
	\]
	
	\textbf{Step 4: Margin of error (99\% confidence, \(z^* = 2.576\)).}
	\[
	E = z^* \cdot SE = 2.576(0.0435) \approx 0.1121
	\]
	
	\textbf{Step 5: Confidence interval.}
	\[
	\hat{p} \pm E = 0.65 \pm 0.1121 \Rightarrow [0.538,\ 0.762]
	\]
	
	\textbf{Interpretation.}
	We are 99\% confident that between about 53.8\% and 76.2\% of households expect inflation to rise over the next year.
	
\end{document}
